% ============================================================
%  related_work.tex — TH-GNN paper, NeurIPS 2025
% ============================================================

\section{Related Work}
\label{sec:related}

% ── Shilling Attack Detection ─────────────────────────────────────────────────
\subsection{Shilling Attack Detection}
\label{sec:rel_shilling}

Shilling attack detection has a two-decade history spanning statistical,
supervised, and graph-based paradigms, surveyed extensively
elsewhere~\cite{si2018shilling,sundar2020understanding,zayed2023experimental,nawara2025shilling}.
The foundational statistical approach treated each user's rating profile as a
feature vector and flagged outliers: Chirita et al.~\cite{chirita2005preventing}
defined per-user features—rating deviation from the item mean, weighted
deviation from mean average (WDMA), profile length, and filler-item
ratio—capturing the coarse fingerprints of injection, PCA-based methods
treated injected profiles as low-rank perturbations separating from genuine
user variance~\cite{williams2006detection}, and clustering or reliability-based
methods grouped users whose rating vectors were suspiciously
similar~\cite{mehta2009unsupervised,zhang2020detecting,cai2021bssc,gunes2025beavar,alonso2019robust}.
These statistical detectors were fast and needed no labelled data but failed
against adaptive attacks calibrated to match the genuine rating distribution,
such as bandwagon~\cite{mobasher2007attacks}.
Supervised classifiers introduced labelled shilling datasets and per-user
feature engineering: random-forest and SVM detectors trained on rating-based
features achieved strong recall against known attack families but overfitted
to the training-time injection signature~\cite{burke2005segment,wen2016svmtia,bansal2021shilldetector}, and
deep-learning variants replaced hand-crafted statistics with recurrent or
convolutional encoders over rating
sequences~\cite{wu2021fight,tong2018shilling,vivekanandan2020hybrid,ebrahimian2020detecting,hao2020unsupervised,hao2018detecting,moradi2023detecting}.
However, all of these approaches treat users independently; they miss the
relational anomalies produced when a group of fake profiles collectively
targets a set of items.

Graph-based detectors exploited this relational signal.
Methods that built user-item bipartite graphs and scored nodes by structural
anomaly measures—degree asymmetry, neighbourhood overlap, co-purchase
concentration—substantially outperformed single-user statistics against
structurally coherent attack families such as AUSH~\cite{lin2020aush}, whose
GAN-based profile generation was subsequently extended by gray-box and
black-box adversarial variants~\cite{wang2022graybox,lin2022shilling,liu2024novel},
penalized latent-perturbation GANs~\cite{nawara2025penalized}, and
multi-distribution generative attacks~\cite{zhou2026mdmgansa,yu2025rupgan}, each
further narrowing the statistical gap between fake and genuine rating
distributions.
On the detection side, group-level and structural methods further exploited
this relational signal via graph convolutional
scoring~\cite{wang2022detecting}, hypergraph spectral
features~\cite{li2021fusing}, and rating-time-series credibility
measures~\cite{zhou2018shilling}.
Anti-FakeU~\cite{zhang2022antifakeu} is the most directly comparable prior
work: it constructs a bipartite user-item graph and detects fake users through
graph-based anomaly scoring over first-order neighbourhood statistics including
co-purchase overlap and degree asymmetry.
Yet Anti-FakeU encodes neither review text nor edge timestamps, leaving two
systematic blind spots: it cannot detect the cross-modal inconsistencies
introduced when LLM-generated text accompanies structurally plausible rating
patterns, and it has no mechanism for recognising the coordinated temporal
bursts that distinguish modern injection campaigns from organic activity.
LLM-based generation has compounded this challenge by eliminating the textual
fingerprints that historically distinguished synthetic from genuine profiles.

% ── LLM-Generated Fake Content Detection ─────────────────────────────────────
\subsection{LLM-Generated Fake Content Detection}
\label{sec:rel_llm_detection}

Detecting computationally generated deceptive content is a longstanding
problem~\cite{ott2011finding,jindal2008opinion} that has grown substantially
harder with the advent of large language models.

The GPT era reset the detection landscape.
Zellers et al.~\cite{zellers2019defending} demonstrated that neural fake-news
generators produce text that fools both human readers and classical detectors,
and that the most effective discriminator is a model of the same family trained
adversarially.
Contemporary LLM-generated text detectors exploit subtle statistical biases
in autoregressive sampling—token-probability curvature, rank-based
features—but these signals vanish under paraphrasing or when the attacker
controls the decoding parameters~\cite{guo2023close}.
LLM watermarking~\cite{kirchenbauer2023watermark} embeds detectable statistical
patterns at generation time by biasing token selection toward a secret hash
partition; while theoretically grounded, watermarks are inapplicable when the
attacker controls the generator and absent from commercially deployed models.

SemanticShield~\cite{li2025semanticshield} targets the recommender-system setting
specifically: it encodes each user's aggregated review corpus with a pre-trained
language model and flags profiles whose representation drifts from the
genuine-user centroid in embedding space.
SemanticShield is the strongest single-modality text baseline for LLM-agent
attacks when review text is present.
However, it degrades to chance on datasets without review text, provides no
mechanism for detecting temporal coordination among injected profiles, and
ignores the structural signals available in the user-item interaction
graph—signals that are informative even when the textual content of fake
profiles is indistinguishable from genuine users.
The structural and temporal graph architectures that can provide these
complementary signals are surveyed next.

% ── Heterogeneous and Temporal GNNs ──────────────────────────────────────────
\subsection{Heterogeneous and Temporal Graph Neural Networks}
\label{sec:rel_gnn}

Standard graph neural networks homogenise all nodes and edges into a single
feature space, discarding semantically meaningful type distinctions.
The Heterogeneous Graph Transformer (HGT; Hu et al.~\cite{hu2020hgt}) addressed
this by parameterising queries, keys, values, and message matrices separately
per node type and per relation type, learning qualitatively different
interaction patterns for each combination.
HGT achieved state-of-the-art performance on heterogeneous academic network
benchmarks and was subsequently adapted for knowledge-graph reasoning and
cross-domain recommendation~\cite{lv2021we}.
The type-specific parameterisation is critical for the shilling detection
setting, where user, item, and review nodes carry structurally and semantically
different attack signatures.

Temporal dynamics introduce a complementary axis.
Temporal Graph Networks (TGN; Rossi et al.~\cite{rossi2020temporal}) maintained
a per-node memory updated by a message function conditioned on
continuous-time edge events, enabling representations that tracked the
evolving state of each node as new interactions arrived.
TGN demonstrated that even simple time-encoding—e.g.\ learnable sinusoidal
embeddings of inter-event gaps—yields large gains on link-prediction benchmarks
with bursty arrival processes.
Fraud detection has independently discovered temporal signals: FRAUDRE and
related approaches showed that the inter-arrival time distribution of fraudulent
accounts differs markedly from genuine users, producing anomaly signals that
static graph features miss entirely~\cite{liu2025fraud}.
Despite these parallel advances, no prior work has combined heterogeneous graph
encoding with learnable temporal attention for shilling attack detection, where
the three-type user-review-item topology and the timestamped injection signal
together provide orthogonal and non-redundant detection axes.
Unlike these works, which use heterogeneous attention as a general-purpose
aggregation mechanism, AAA treats anomaly detection as a first-class objective
within the attention computation—encoding domain-specific signals directly
into the attention score rather than as post-hoc feature streams.

% ── LLM Agents for Recommendation Attack ─────────────────────────────────────
\subsection{LLM Agents for Recommendation Attack}
\label{sec:rel_llm_attack}

Agent4SR~\cite{zhang2024agent4sr} was the first work to frame shilling attack
construction as a multi-step agentic task: a large language model was prompted
to construct a complete fake user persona, generate a coherent review history,
and assign ratings consistent with the persona's stated preferences, all while
optimising target-item exposure.
The resulting profiles were grammatically fluent, topically coherent, and
statistically similar to genuine users, systematically evading detectors
that relied on either distributional rating anomalies or linguistic
outlier detection.
This capability builds on the broader emergence of LLM-based and agentic
recommendation~\cite{lin2023survey,wang2023survey,plaat2025agentic}, which has
made such attacks increasingly practical to mount at scale, and raises
security concerns that extend beyond recommender systems to LLM agents
generally~\cite{he2024security}.
Related work has explored LLM-based poisoning in broader adversarial
recommendation settings~\cite{deldjoo2024survey}, and prompt-injection attacks
against retrieval-augmented recommenders~\cite{shi2023large}, but systematic
evaluation of detection strategies against agent-generated profiles across
multiple datasets and injection rates remains absent from the literature.

% ── Gap synthesis ─────────────────────────────────────────────────────────────

The gaps surveyed above form three orthogonal failure modes in existing
detectors: Anti-FakeU's reliance on graph structure alone leaves it blind to
the cross-modal inconsistencies introduced by LLM-generated review text;
SemanticShield's reliance on text alone renders it ineffective on rating-only
datasets and unable to leverage graph topology; and neither family incorporates
temporal information, missing the coordinated injection-timing signals that are
present even when individual profiles appear genuine.
TH-GNN closes all three gaps simultaneously: the two-layer HGT backbone encodes
the heterogeneous user-review-item graph with per-type and per-relation
parameterisation, the cross-modal attention fusion integrates review and item
text streams with the structural embedding, and the sinusoidal temporal
encoding together with the burstiness GRU capture the edge-timestamp patterns
that betray injection campaigns at both the per-edge and per-user level.
The result is a detector that is strictly more informative than any
single-modality baseline and, as the ablation study confirms, degrades
gracefully when any individual stream is absent rather than failing
catastrophically.
